\documentclass{bmvc2k}

\usepackage{graphicx}
\usepackage{booktabs}

\title{From Embeddings to Recurrent Patterns:\\
A Research Portfolio Note on Unsupervised Chat Clustering}

\addauthor{Bc. Mykola Vorontsov}{m.vorontsov@campus.tu-berlin.de}{1}
\addinstitution{TU Berlin\\Berlin-Charlottenburg, Germany}

\runninghead{Vorontsov}{Embeddings to Recurrent Patterns}

\def\eg{\emph{e.g}\bmvaOneDot}
\def\Eg{\emph{E.g}\bmvaOneDot}
\def\etal{\emph{et al}\bmvaOneDot}

\begin{document}

\maketitle

\begin{abstract}
This note documents a portfolio-style machine learning project on unsupervised discovery of recurring phrase-like and topic-like patterns in a private Russian Telegram group chat. The practical motivation is to understand whether short informal messages form dense semantic neighborhoods in sentence-embedding space, and whether those neighborhoods can recover repeated conversational expressions beyond broad topics. We compare K-Means, HDBSCAN on raw embeddings, UMAP followed by HDBSCAN, BERTopic, and DBSCAN with explicit radius control. The main result is qualitative but consistent: density-based methods become substantially more useful after dimensionality reduction, and DBSCAN provides the clearest control over the trade-off between broad thematic clusters and narrow recurring expressions. At the same time, the project remains deliberately honest about its limits: the dataset is private, the evaluation is partly qualitative, and temporal trend detection is not yet solved well enough to be claimed as a finished contribution.
\end{abstract}

\section{Introduction}
\label{sec:intro}

Informal group chats contain many recurring linguistic behaviors: exact repeated phrases, lightly varied catchphrases, narrow in-group templates, and also many generic conversational fillers. The distinction matters. A useful system for recurring-content discovery should surface the first two categories more reliably than the last one.

This repository studies that problem as an unsupervised clustering task on sentence embeddings. The aim is not to present a finished conference paper, but to document a careful research workflow and its current conclusions. The framing was therefore narrowed from ``meme/trend detection'' to a more defensible goal: \emph{discovery of recurring phrase-like and topic-like patterns in short chat messages}. This narrower formulation matches what the current experiments can genuinely support.

The project uses a private Telegram export, multilingual sentence embeddings, and several clustering pipelines. The central working hypothesis is that repeated conversational patterns occupy locally dense regions in embedding space, while unrelated short messages appear as sparse noise. Density-based clustering is therefore a natural fit, but only if the geometry of the embedding space is suitable for neighborhood-based reasoning.

The main contributions of the current repository version are:
\begin{itemize}
    \item a clean and reproducible experiment pipeline for comparing baseline and density-based clustering methods;
    \item a public, code-only portfolio packaging that avoids releasing private chat data;
    \item a more honest problem statement aligned with what the experiments actually demonstrate;
    \item a lightweight qualitative evaluation protocol for auditing whether recovered clusters are phrase-like, topic-like, generic conversational, or artifact-heavy.
\end{itemize}

\section{Related Work}

The project sits at the intersection of sentence embeddings, density-based clustering, and neural topic modeling. Sentence-BERT-style embeddings provide a practical way to map semantically related short utterances into a shared vector space~\cite{reimers2019sentencebert}. This is particularly useful for chat data, where lexical overlap alone is often too brittle because repeated expressions may appear with spelling variation, slang, or paraphrase.

For clustering, the comparison focuses on methods with different inductive biases. K-Means offers a simple baseline, but it assumes approximately spherical clusters and requires a predefined number of clusters. DBSCAN~\cite{ester1996dbscan} and HDBSCAN~\cite{campello2015hdbscan} instead define clusters through local density, making them better suited to irregular neighborhoods and noise-heavy datasets. In practice, however, density estimation can become unstable in high-dimensional embedding spaces.

UMAP~\cite{mcinnes2018umap} is therefore used as a geometry-shaping step before density-based clustering. Its role here is pragmatic rather than theoretical: preserving enough local structure that repeated conversational patterns become easier to isolate. BERTopic~\cite{grootendorst2022bertopic} is included not because the task is classical topic modeling, but because it provides a convenient integrated workflow around embeddings, dimensionality reduction, clustering, and topic representation.

This project differs from standard topic-modeling use cases in one important way: the target phenomena are often narrower than topics. Some clusters correspond to recognizable discussion themes, but the more interesting cases are repeated phrase templates and conversational catchphrases. That is precisely why a qualitative audit remains necessary.

\section{Data and Public Release Constraints}
\label{sec:data}

The underlying dataset consists of messages extracted from a private Telegram group chat with 19 distinct authors visible in the processed export over a period of six months. The messages are predominantly in Russian. The original export contains 42{,}698 entries. After filtering to text messages and removing empty or very short content, 32{,}738 messages remain for analysis.

Because the source material is private, the public repository is intentionally code-only. It does not include:
\begin{itemize}
    \item the Telegram export itself,
    \item user identifiers,
    \item media files,
    \item raw message quotations.
\end{itemize}

Public artifacts therefore contain only aggregate metrics, paraphrased illustrative descriptions, and documentation of the experiment pipeline. This is a core design decision of the repository, not a temporary omission.

\section{Methodology}

\subsection{Preprocessing and Embeddings}

Messages are filtered to entries with type \texttt{message}, converted into plain text, and trimmed to a minimum cleaned length of six characters. This threshold is a pragmatic compromise: it removes a large amount of trivial noise, but it also risks discarding some genuine short recurring expressions. The note therefore treats this preprocessing choice as a limitation rather than as an optimal design.

Sentence embeddings are computed with the multilingual model \texttt{paraphrase-multilingual-MiniLM-L12-v2}. The model is chosen because it supports Russian and provides a practical semantic representation for short informal messages.

\subsection{Clustering Pipelines}

The canonical experiment pipeline contains four entrypoints in the repository CLI:
\begin{enumerate}
    \item \texttt{preprocess}: summarize filtering and dataset size;
    \item \texttt{baseline}: run K-Means baselines;
    \item \texttt{density}: run HDBSCAN, UMAP + HDBSCAN, BERTopic, and DBSCAN sweeps;
    \item \texttt{report}: export the full private summary bundle.
\end{enumerate}

The compared methods are:
\begin{itemize}
    \item K-Means as a simple semantic baseline;
    \item HDBSCAN on raw embeddings;
    \item UMAP followed by HDBSCAN;
    \item BERTopic with its default HDBSCAN-based clustering flow;
    \item DBSCAN after UMAP with several values of $\epsilon$ and fixed \texttt{min\_samples=3}.
\end{itemize}

\subsection{Qualitative Evaluation Protocol}

Quantitative summaries alone are not sufficient for this task. A clustering method can produce compact clusters and still fail the actual goal if those clusters correspond mainly to generic conversational filler or fragmented artifacts.

To address that, the public repository includes a small manual audit protocol. A sampled cluster is assigned one of four labels:
\begin{itemize}
    \item \texttt{phrase-like},
    \item \texttt{topic-like},
    \item \texttt{generic conversational},
    \item \texttt{artifact/noise}.
\end{itemize}

The purpose of this audit is not benchmark-grade evaluation. Its purpose is narrower: to inspect whether a method surfaces clusters that are plausibly useful for recurring phrase discovery.

\section{Results}

\subsection{Baseline}

K-Means is easy to interpret as a first pass, but it mostly recovers broad semantic regions rather than tight repeated expressions. For $K=100$, the median cluster size is 264 messages and the mean cluster size is 327.38, which is far too coarse for the phrase-discovery goal. Larger values of $K$ fragment the dataset more aggressively, but still do not provide a principled notion of noise or density.

\subsection{Density-Based Methods}

HDBSCAN on raw embeddings is too conservative for this dataset. It labels 28{,}631 of 32{,}738 messages as noise, while the median non-noise cluster size is only 7. This suggests that the original 384-dimensional embedding space is too sparse for stable density estimation.

Applying UMAP before HDBSCAN changes the picture considerably. Noise drops to 9{,}961 messages and the method discovers 1{,}378 clusters with median size 10. The improvement is substantial, but fragmentation remains visible: many narrow conversational patterns are still split into separate micro-clusters.

BERTopic behaves similarly to UMAP + HDBSCAN in this setting. In the current run it produces 1{,}438 clusters, labels 10{,}549 messages as noise, and keeps the median non-noise cluster size at 10. Its main advantage is convenience and inspectability rather than a strong change in clustering quality. It is therefore helpful as an analysis wrapper, but not a decisive modeling improvement by itself.

DBSCAN after UMAP gives the most useful control over cluster granularity. The effect of $\epsilon$ is summarized in Table~\ref{tab:dbscan}.

\begin{table}[h]
\centering
\begin{tabular}{lccc}
\toprule
$\epsilon$ & \# Clusters & Median Size & Main Behavior \\
\midrule
0.01  & 953 & 5 & more localized, still partly thematic \\
0.005 & 602 & 4 & best trade-off for phrase discovery \\
0.001 & 175 & 5 & overly strict, dominated by trivial reactions \\
\bottomrule
\end{tabular}
\caption{DBSCAN behavior after UMAP dimensionality reduction.}
\label{tab:dbscan}
\end{table}

The interpretation is straightforward: larger neighborhoods keep more thematic content together, while smaller neighborhoods isolate only the most compact repeated expressions. In the current experiments, $\epsilon = 0.005$ is the most convincing balance between specificity and interpretability.

\section{Discussion}

The strongest repository-level conclusion is not that one method ``solves'' recurring content discovery, but that the geometry of the representation space determines whether density-based clustering is even viable. Raw sentence embeddings are too sparse for this task. UMAP helps. DBSCAN then exposes a practically useful granularity knob that HDBSCAN alone does not provide as clearly.

At the same time, the experiments also show a more sobering point: many recovered clusters are neither canonical topics nor clean meme phrases. They often sit somewhere in between. That ambiguity is not a weakness of documentation; it is one of the main empirical findings of the project.

\section{Limitations and Ethics}

This project has four major limitations.

First, the dataset is private and small in the sense of conversational diversity, even if it is reasonably large in message count. Claims should therefore be read as a case study, not as a general benchmark conclusion.

Second, the evaluation remains partly qualitative. The repository improves this with a transparent audit protocol, but it still does not replace a large-scale labeled benchmark.

Third, the current preprocessing rule removes short messages below six characters. This likely discards some real recurring signals such as short laughter variants or reaction phrases.

Fourth, temporal trend modeling is intentionally left as future work. The current repository can identify recurring local structure, but it does not yet provide a robust notion of emergence, decay, or burstiness over time.

On the ethics side, the public release is designed to avoid exposing private conversation content. No raw export, user identifiers, or direct quotations are needed to understand the method or its limitations.

\section{Reproducibility}

The root repository now exposes a single canonical CLI pipeline and a code-only public artifact set. A user with local access to the private Telegram export can reproduce the experiment summaries via:
\begin{verbatim}
python main.py preprocess
python main.py baseline
python main.py density
python main.py report
\end{verbatim}

The public-facing reader, by contrast, can inspect the repository through the sanitized summary files and the portfolio note without needing access to any private chat data.

\bibliography{egbib}
\end{document}
